\section{Discussion}
\label{sec:discussion}

\subsection{What Our Results Do and Do Not Show}

% Important: this is the constructive reframing per docs/paper_strategy.md.

It is tempting to read our results as ``foundation models cannot detect
stress from resting-state EEG.'' That framing is incorrect, and the
distinction matters.

\paragraph{What we do show.} Two completely different methods --- a
200M-parameter pretrained transformer and a random forest on hand-crafted
band-power features --- converge to the same balanced accuracy
($\sim$0.66) on subject-level cross-validation with $N{=}17$. If no
diagnostic signal existed, both would collapse to chance ($\sim$0.50).
The fact that they converge to a number above chance, and to
\emph{the same} number, is positive evidence that a real cross-subject
resting-state stress signature exists in our data. The ceiling is
\emph{method-independent}, not method-failure.

\paragraph{What we do not show.} We do not show that the underlying
biology is silent, that EEG is the wrong modality, or that the dream of
pre-diagnosing prolonged stress is unrealistic. We show only that the
specific framing --- binary, single-recording, generic-feature,
cross-subject zero-shot prediction with $N{=}17$ --- saturates near
$0.66$ regardless of model.

\subsection{Why Subject Identity Dominates}

Foundation models pretrained on masked patch prediction or contrastive
objectives are rewarded for any feature that helps reconstruct or
discriminate windows. Subject-specific traits (skull thickness, electrode
impedance, individual spectral fingerprints) are stable and discriminative;
clinical labels are less so. The pretraining objective therefore implicitly
selects for subject identity --- and the bigger the model, the more
faithfully it encodes it.

This is the same phenomenon documented in rsfMRI ``brain
fingerprinting''~\citep{finn2015}: individual functional connectomes are
remarkably stable across sessions, and accurately identifiable across
hundreds of subjects, while task-related variance is comparatively small.

\subsection{Why Standard Mitigations Fail}

Subject-adversarial methods (GRL, LEAD-style auxiliary losses) attempt to
remove subject information from the learned features. Our results
(Table~\ref{tab:adv}) show this hurts rather than helps.

We hypothesise the explanation is straightforward: removing variance does
not create signal. If $\eta^{2}_{\mathrm{label}} \approx 0.06$ in the raw
features, then no projection of those features can exceed that label
information content. The subject information is acting as a stable
\emph{reference frame}, and removing it leaves the small label variance
exposed but unamplified --- and now also less robust to noise.

\subsection{Implications: Reframe the Question}

The constructive read of our findings is that subject identity dominance is
a structural feature, not a bug, and the right response is to change the
question being asked of the model:

\begin{itemize}[topsep=0.25em,itemsep=0.1em,leftmargin=1.5em]
  \item \textbf{Within-subject longitudinal modelling.} Use each subject as
    their own reference. Predict the deviation of a follow-up recording
    from the subject's baseline rather than absolute label. This converts
    subject identity from noise into the reference frame and matches the
    actual clinical use case (pre-diagnose stress in \emph{this individual}
    over time). We list this as the highest-leverage open experiment in
    Section~\ref{sec:future}.
  \item \textbf{Personalised few-shot adaptation.} Standard in clinical BCIs
    but underused for psychiatric EEG. A handful of calibration samples per
    subject may unlock the per-subject signal that cross-subject models
    cannot reach.
  \item \textbf{Continuous label regression.} Binary thresholding throws
    away most of the label variance. Regressing against the actual DASS
    score is a natural alternative.
  \item \textbf{Multimodal fusion.} EEG's modest signal may be detectable
    when combined with autonomic markers (HRV, actigraphy). The diagnosis
    here suggests EEG should be one input among several, not the sole
    input.
\end{itemize}

\subsection{Limitations}
\label{sec:limitations}

\begin{itemize}[topsep=0.25em,itemsep=0.1em,leftmargin=1.5em]
  \item Small subject count ($N{=}17$) on the primary dataset. We address
    this by validating the variance pattern on two larger reference
    datasets (ADFTD $N{=}65$, TDBRAIN $N{=}359$).
  \item Three FMs is not exhaustive. We chose architectures spanning ViT
    (LaBraM, REVE) and criss-cross transformers (CBraMod), but recent
    models such as BIOT and LARGE-Brain are not evaluated here.
  \item The 19-channel common montage limits frozen-feature analysis to a
    subset of available electrodes. Per-dataset fine-tuning uses the full
    montage.
  \item We only evaluate full fine-tuning. Linear probing and adapter
    methods may behave differently and are an obvious extension.
\end{itemize}

\subsection{Future Work}
\label{sec:future}

% This section is the bridge to the next paper (or the next iteration of
% this one). Do not delete the within-subject longitudinal proposal --- it
% is the main constructive output of the diagnosis.

\begin{enumerate}[topsep=0.25em,itemsep=0.1em,leftmargin=1.5em]
  \item \textbf{Within-subject longitudinal modelling on UCSD Stress
    (highest leverage).} Compute per-subject baseline features, predict
    deviation from baseline, evaluate under the same subject-level CV
    protocol. If \BA{} exceeds 0.72, the productive next paper writes
    itself.
  \item \textbf{Spectral-guided FiLM conditioning.} Inject neurophysiology
    priors (theta/alpha/beta band power) into the FM via FiLM layers.
    Whether this works is an open question.
  \item \textbf{Subject-aware contrastive pretraining on small clinical
    cohorts.} Flip the diagnosis on its head: explicitly encode subject as
    a feature and condition the classifier on it.
  \item \textbf{Replication on additional resting-state stress cohorts}
    (EEGMAT, SAM40, EDPMSC), even though these use acute stress
    induction --- to test whether the ceiling is universal or
    chronic-stress-specific.
\end{enumerate}
